\section{Introduction}

% ============================================================
% No explicit word limit, but keep it short — the marks are in
% sections 1 and 2. 1-2 paragraphs is plenty.
% Excluded from word count: bibliography, equations, figures.
% ============================================================

% SUGGESTED CONTENT:
%
%   Para 1 — Context and motivation (~3 sentences)
%     Describe the domain gap between game and movie footage.
%     Explain why human regions are the focus (foreground style matters
%     most; background mixing hurts quality in full-frame translation).
%     One sentence on the broader application (game realism enhancement).
%
%   Para 2 — Approach summary (~3 sentences)
%     "We present a pipeline that: (1) extracts human patches via
%      YOLOv8~\cite{jocher2023} with quality-aware selection; (2)
%      classifies patches into five pose categories using a GCN trained
%      on COCO keypoints; (3) selects diverse representative patches via
%      DINOv2 clustering; and (4) applies CUT~\cite{park2020}
%      patch-level style transfer with EMA temporal blending."
%     Reference Figure showing the overall pipeline if you make one.
%
% KEY REFERENCES TO CITE:
%   - CUT: Park et al. (2020), "Contrastive Unpaired Translation" — the core model
%   - CycleGAN: Zhu et al. (2017) — precursor, good for context
%   - YOLOv8: Jocher et al. (2023), ultralytics/ultralytics
%   - DINOv2: Oquab et al. (2023) — for 1.3 embedding
%   - GCN: Kipf & Welling (2017), "Semi-Supervised Classification with GCNs"
%   - ST-GCN (if you implement temporal GCN): Yan et al. (2018)
%
% OPTIONAL FIGURE (no word cost):
%   A pipeline diagram showing the flow 1.1→1.2→1.3→2.1→2.2.
%   Save to: paper/figs/pipeline_overview.pdf
